\documentclass[a4paper]{article}

\usepackage{graphicx}
\usepackage{amsmath,amssymb}
\usepackage{minted}
\usepackage{multirow}
\usepackage{float}
\usepackage{todonotes}
\usepackage{forest}
\usepackage{xstring}
\usepackage{xcolor}
\usepackage[most]{tcolorbox}
\usepackage{xparse}
\usepackage{natbib}

\usetikzlibrary{graphs,graphdrawing}
\usegdlibrary{layered}

% for external graphviz
\input{/home/donkarlo/Dropbox/repo/nd_language_ai_project/src/nd_language_ai/formal/computer/programming/tex/latex/code_clips/external_graphviz.tex}

% for tags
\input{/home/donkarlo/Dropbox/repo/nd_language_ai_project/src/nd_language_ai/formal/computer/programming/tex/latex/part/control_sequence/kind/semtag/semtag.tex}

% for links
\input{/home/donkarlo/Dropbox/repo/nd_language_ai_project/src/nd_language_ai/formal/computer/programming/tex/latex/code_clips/seehere.tex}

\title{}
\date{}

\begin{document}
   \maketitle

   The grandmother-cell hypothesis proposes that highly selective neurons can represent specific complex objects or concepts.
   Gross reviews the historical development of this idea and distinguishes it from more realistic sparse neuronal representations.
   The paper argues that highly selective neurons may participate in recognition without requiring exactly one neuron for each concept \cite{gross-2002-genealogy-of-the-grandmother-cell}.

   Single neurons in the human medial temporal lobe were found to respond selectively to particular individuals, objects, or landmarks.
   Some neurons responded invariantly to different images or representations of the same concept, such as the well-known ``Jennifer Aniston neuron.''
   These findings support sparse and abstract neuronal representations of concepts, closely related to the grandmother-cell hypothesis \cite{quiroga-2005-invariant-visual-representation-by-single-neurons-in-the-human-brain}.

            \paragraph{Wikipedia}
               is a hypothetical neuron that represents a complex but specific concept or object \seeherefooter{https://en.wikipedia.org/wiki/Grandmother_cell}


               \bibliographystyle{plainnat}
               \bibliography{/home/donkarlo/Dropbox/repo/nd_shared_research/bibliography/bibliography.bib}
\end{document}